\section*{अध्याय \ref{ch:g10:universal-dna} --- DNA: एक सार्वत्रिक आनुवंशिक अणु}

\begin{solution}{exo:g10:universal-dna:1}
एक फ़ॉस्फ़ेट समूह, एक शर्करा (डिऑक्सीराइबोस) और एक \omterm{def:g10:universal-dna:nucleotide}{क्षारक}। \omterm{def:g10:universal-dna:nucleotide}{क्षारक}:
एडेनीन, थाइमीन, ग्वानीन, साइटोसीन।
\end{solution}

\begin{solution}{exo:g10:universal-dna:2}
\texttt{ATGCCTAAGT}, \omterm{def:g10:universal-dna:nucleotide}{क्षारक} के नीचे \omterm{def:g10:universal-dna:nucleotide}{क्षारक} लिखते हुए।
\end{solution}

\begin{solution}{exo:g10:universal-dna:3}
C $= 21\%$; A $=$ T $= (100 - 42)/2 = 29\%$ हर एक।
\end{solution}

\begin{solution}{exo:g10:universal-dna:4}
\omterm{def:g10:universal-dna:gene}{जीन} \omterm{def:g10:universal-dna:information}{DNA} का वह खंड है जिसका अनुक्रम एक \omterm{def:g10:chemistry-of-life:families}{प्रोटीन} के निर्देश ढोता है और जो
किसी \omterm{def:g10:universal-dna:information}{गुणसूत्र} पर एक तय स्थान पर बैठा होता है। \omterm{def:g10:universal-dna:gene}{युग्मविकल्पी} किसी \omterm{def:g10:universal-dna:gene}{जीन} के उन
रूपों में से एक है जो अपने अनुक्रम के चंद \omterm{def:g10:universal-dna:nucleotide}{क्षारकों} से बाक़ी रूपों से अलग
पड़ता है।
\end{solution}

\begin{solution}{exo:g10:universal-dna:5}
वह जीव जिसके \omterm{def:g10:universal-dna:information}{DNA} में किसी दूसरी जाति का \omterm{def:g10:universal-dna:gene}{जीन} डाला गया हो और जो उसे पढ़ता
तथा आगे पहुँचाता हो। उदाहरण: जेलीफ़िश का GFP \omterm{def:g10:universal-dna:gene}{जीन} ढोते चूहे; मानव इंसुलिन
का \omterm{def:g10:universal-dna:gene}{जीन} ढोता \emph{E.~coli}।
\end{solution}

\begin{solution}{exo:g10:universal-dna:6}
$1500 \times 0.34 = \qty{510}{nm}$; $1500/10 = 150$ चक्कर।
\end{solution}

\begin{solution}{exo:g10:universal-dna:7}
$2 \times \num{3.2e9} \times \qty{0.34}{nm} \approx \qty{2.2}{m}$, यानी
\omterm{def:g10:cells-common-unit:organelle}{केंद्रक} के व्यास का लगभग $\num{3.6e5}$ गुना।
\end{solution}

\begin{solution}{exo:g10:universal-dna:8}
किसी ``शुद्ध किए गए'' अंश में दूसरे अणुओं के अंश अब भी हो सकते हैं,
इसलिए अकेली उसकी सक्रियता यह सिद्ध नहीं करती कि ज़िम्मेदार \omterm{def:g10:universal-dna:information}{DNA} ही है।
पूरे सत्त में ख़ास तौर पर \omterm{def:g10:universal-dna:information}{DNA} को नष्ट करके यह देखना कि सक्रियता ग़ायब हो
गई, जबकि \omterm{def:g10:chemistry-of-life:families}{प्रोटीन} और RNA नष्ट करने वाले एंजाइम उसे छोड़ देते हैं, असर को
\omterm{def:g10:universal-dna:information}{DNA} पर और किसी और पर नहीं टिका देता है।
\end{solution}

\begin{solution}{exo:g10:universal-dna:9}
कोई नियत दोहराव हर जाति में एक ही संघटन (हर \omterm{def:g10:universal-dna:nucleotide}{क्षारक} का 25\%) देता।
शारगाफ़ ने A+T को जातियों के बीच काफ़ी फैला हुआ पाया: \omterm{def:g10:universal-dna:nucleotide}{क्षारकों} का क्रम
कोई नियत ढर्रा नहीं है, इसलिए अनुक्रम बदलने को स्वतंत्र है --- और यही
उसे जानकारी ढोने के लायक़ बनाता है।
\end{solution}

\begin{solution}{exo:g10:universal-dna:10}
\omterm{def:g10:chemistry-of-life:families}{प्रोटीन} बनाने के लिए अनुक्रम क्रम से पढ़ा जाता है; जिस स्थान पर कोई
\omterm{def:g10:chemistry-of-life:families}{ऐमीनो अम्ल} तय होता है वहाँ बदला हुआ \omterm{def:g10:universal-dna:nucleotide}{क्षारक} एक \omterm{def:g10:chemistry-of-life:families}{ऐमीनो अम्ल} की जगह दूसरा
बिठा सकता है और \omterm{def:g10:chemistry-of-life:families}{प्रोटीन} का आकार या सक्रियता बदल सकता है। पर यदि बदला हुआ
\omterm{def:g10:universal-dna:nucleotide}{क्षारक} \omterm{def:g10:universal-dna:gene}{जीन} के किसी न पढ़े जाने वाले हिस्से में हो, या वही \omterm{def:g10:chemistry-of-life:families}{ऐमीनो अम्ल} तय
करता हो, तो \omterm{def:g10:chemistry-of-life:families}{प्रोटीन} अनबदला रहता है।
\end{solution}

\begin{solution}{exo:g10:universal-dna:11}
चूहे की सारी \omterm{def:g10:cells-common-unit:cell}{कोशिकाएँ} निषेचित अंडाणु से विभाजन द्वारा उतरी हैं; हर
विभाजन पर \omterm{def:g10:universal-dna:information}{DNA} की, पारजीन समेत, नक़ल (\omterm{prop:g10:universal-dna:helix}{पूरकता} से) बनती है और वह बँट जाता
है। इसलिए त्वचा, यकृत और शुक्राणु --- सब उस \omterm{def:g10:universal-dna:gene}{जीन} को ढोते हैं, और शुक्राणु
उसे अगली पीढ़ी तक पहुँचा देता है।
\end{solution}

\begin{solution}{exo:g10:universal-dna:12}
$\num{4.6e6} \times \qty{0.34}{nm} \approx \qty{1.6}{mm}$, यानी \omterm{def:g10:cells-common-unit:cell}{कोशिका}
की लंबाई का क़रीब 800 गुना। पर अणु केवल \qty{2}{nm} चौड़ा है: उसका
आयतन, लगभग $\pi \times 1^2 \times \num{1.6e6}\,
\unit{nm^3} \approx \num{5e-3}\,\unit{\micro m^3}$, \omterm{def:g10:cells-common-unit:cell}{कोशिका} के आयतन का
क़रीब 1\% है; मोड़ा और कुंडलित किया हुआ वह जगह बचाकर समा जाता है।
\end{solution}

\begin{solution}{exo:g10:universal-dna:13}
मानव \omterm{def:g10:universal-dna:information}{DNA} 700 गुना लंबा है पर उसमें \omterm{def:g10:universal-dna:gene}{जीन} केवल लगभग 5 गुना अधिक हैं: उसका
अधिकांश हिस्सा \omterm{def:g10:universal-dna:gene}{जीन} है ही नहीं। उस अतिरिक्त \omterm{def:g10:universal-dna:information}{DNA} का एक हिस्सा यह नियंत्रित
करता है कि \omterm{def:g10:universal-dna:gene}{जीन} कब इस्तेमाल हों, और एक हिस्सा दोहराया हुआ अनुक्रम है
जिसका कोई ज्ञात काम नहीं। \omterm{def:g10:universal-dna:nucleotide}{क्षारकों} की गिनती इबारत का आकार नापती है, जीव
की जटिलता नहीं; \omterm{def:g10:universal-dna:gene}{जीनों} की संख्या और सबसे बढ़कर उनके मेल तथा नियंत्रण का
ढंग कहीं अधिक मायने रखते हैं।
\end{solution}

\begin{solution}{exo:g10:universal-dna:14}
वह \omterm{def:g10:universal-dna:information}{DNA} की सार्वत्रिकता का खंडन करती
(\cref{prop:g10:universal-dna:universal})। चूँकि साझा अणु और साझा कूट
किसी साझा पूर्वज से विरासत में मिलते हैं, इसलिए ऐसी जाति को जीवन के किसी
अलग उद्गम से उतरा होना पड़ता: उसके \omterm{def:g10:universal-dna:gene}{जीन} हमारी \omterm{def:g10:cells-common-unit:cell}{कोशिकाएँ} पढ़ ही नहीं
पातीं, और दोनों वंशक्रमों के बीच \omterm{prop:g10:universal-dna:universal}{पारजीनन} असंभव होता।
\end{solution}

\begin{solution}{exo:g10:universal-dna:15}
जीवाणु वाला उत्पाद ठीक-ठीक वही मानव \omterm{def:g10:chemistry-of-life:families}{प्रोटीन} है (सूअर का इंसुलिन एक \omterm{def:g10:chemistry-of-life:families}{ऐमीनो
अम्ल} से अलग होता है और प्रतिरक्षा प्रतिक्रियाएँ भड़का सकता है), और वह
असीमित मात्रा में तथा सूअर के रोगजनकों से मुक्त बनता है। कोई नियामक फिर
भी पूछता कि शुद्ध किया गया \omterm{def:g10:chemistry-of-life:families}{प्रोटीन} जीवाणु के संदूषकों से मुक्त है या
नहीं, और वह ठीक से मुड़ा हुआ तथा सक्रिय है या नहीं।
\end{solution}

\begin{solution}{pb:g10:universal-dna:1}
\textbf{1.} $720 \times 0.34 \approx \qty{245}{nm}$; $720/10 = 72$
चक्कर।

\textbf{2.} \texttt{TACTCATTTCCTCTTCTTG}।

\textbf{3.} T $= 27\%$; G $=$ C $= (100 - 54)/2 = 23\%$।

\textbf{4.} $720/238 \approx 3.0$: यानी प्रति \omterm{def:g10:chemistry-of-life:families}{ऐमीनो अम्ल} तीन \omterm{def:g10:universal-dna:nucleotide}{क्षारक}
(आख़िरी तीन शृंखला का अंत बताते हैं)।

\textbf{5.} $4^{720}$ अनुक्रम, यानी लगभग $10^{433}$। रसायन किसी भी क्रम
की छूट देता है, इसलिए इतनी सब संभावनाओं में से चुना गया वह ख़ास क्रम एक
संदेश है --- यानी जानकारी।

\textbf{6.} सार्वत्रिकता: हर सजीव उन्हीं चार \omterm{def:g10:universal-dna:nucleotide}{न्यूक्लियोटाइडों} का
इस्तेमाल करता है और किसी अनुक्रम को उन्हीं नियमों से पढ़ता है।

\textbf{7.} चूहे की सारी \omterm{def:g10:cells-common-unit:cell}{कोशिकाएँ} अंडाणु से विभाजन द्वारा उतरी हैं,
इसलिए अंडाणु में डाला गया \omterm{def:g10:universal-dna:gene}{जीन} हर \omterm{def:g10:cells-common-unit:cell}{कोशिका} में नक़ल हो जाता है। वयस्क की
त्वचा में डाला गया वह केवल उपचारित त्वचा \omterm{def:g10:cells-common-unit:cell}{कोशिकाओं} और उनके वंशजों में ही
होता।

\textbf{8.} हर तरह की \omterm{def:g10:cells-common-unit:cell}{कोशिका} में पूरा \omterm{def:g10:universal-dna:information}{DNA} रहता है, पारजीन समेत:
विशेषीकरण \omterm{def:g10:universal-dna:gene}{जीनों} को फेंकता नहीं।

\textbf{9.} हाँ: जनन \omterm{def:g10:cells-common-unit:cell}{कोशिकाएँ} वह \omterm{def:g10:universal-dna:gene}{जीन} ढोती हैं। यदि चूहे में जोड़े के एक
\omterm{def:g10:universal-dna:information}{गुणसूत्र} पर उसकी एक प्रति हो, तो उसकी लगभग आधी संतानें उसे पाएँगी और
चमकेंगी।

\textbf{10.} \omterm{def:g10:chemistry-of-life:families}{प्रोटीन} जानकारी नहीं है: वह ख़र्च हो जाता है और दोबारा कभी
नहीं बनता, इसलिए चमक फीकी पड़ जाती है। \omterm{def:g10:universal-dna:gene}{जीन} की नक़ल हर विभाजन पर बनती है
और वह लगातार पढ़ा जाता है, इसलिए चमक स्थायी और वंशागत होती है।

\textbf{11.} \omterm{def:g10:universal-dna:information}{DNA} वाला अंश; और \omterm{def:g10:universal-dna:information}{DNA} को नष्ट करने वाले एंजाइम से उपचार।

\textbf{12.} जीवाणुओं ने दूसरे प्रभेद का \omterm{def:g10:universal-dna:information}{DNA} ग्रहण किया, उसे अपने में
जोड़ा, उस लक्षण को व्यक्त किया और आगे पहुँचाया: यानी एक जीव से दूसरे में
\omterm{def:g10:universal-dna:gene}{जीन} का अंतरण, और \omterm{prop:g10:universal-dna:universal}{पारजीनन} जानबूझकर यही करता है।

\textbf{13.} A$\,=\,$T और G$\,=\,$C क्षारक-जोड़ी का समर्थन करते हैं;
बदलता हुआ A+T हिस्सा दिखाता है कि जोड़ों का क्रम स्वतंत्र है।

\textbf{14.} A और G दो-वलय वाले बड़े \omterm{def:g10:universal-dna:nucleotide}{क्षारक} हैं, T और C एक-वलय वाले
छोटे: कोई बड़ा किसी छोटे के साथ जुड़े तो दूरी हमेशा एक ही बनती है, इसलिए
दोनों रीढ़ें पूरे अणु भर \qty{2}{nm} की दूरी पर बनी रहती हैं। दो बड़े
\omterm{def:g10:universal-dna:nucleotide}{क्षारक} उभार बनाते और दो छोटे एक ख़ाली जगह छोड़ देते।

\textbf{15.} क्योंकि हर लड़ी दूसरी की पूरक है: उन्हें अलग कीजिए, जोड़ी
के नियम से हर एक पर साथी बनाइए, और दो एक जैसे अणु बन जाते हैं।

\textbf{16.} $2 \times \num{2.7e9} \times \qty{0.34}{nm} \approx
\qty{1.8}{m}$।

\textbf{17.} $\qty{1.8}{m} / \qty{6}{\micro m} = \num{3e5}$: यानी तीन
लाख गुना सघनन।

\textbf{18.} $720 / \num{5.4e9} \approx \num{1.3e-7}$: यानी \omterm{def:g10:cells-common-unit:cell}{कोशिका} के
\omterm{def:g10:universal-dna:information}{DNA} का लगभग एक करोड़वाँ हिस्सा।

\textbf{19.} $\num{2e4} \times \num{3e4} = \num{6e8}$ क्षारक-जोड़े,
यानी जीनोम का लगभग 22\%; बाक़ी \omterm{def:g10:universal-dna:gene}{जीनों} के बीच में पड़ा है।

\textbf{20.} चूहे की एक \omterm{def:g10:cells-common-unit:cell}{कोशिका} लगभग \qty{1.8}{m} \omterm{def:g10:universal-dna:information}{DNA} ढोती है; और चूँकि
वह \omterm{def:g10:universal-dna:information}{DNA} ठीक वैसे ही बना और वैसे ही पढ़ा जाता है जैसे जेलीफ़िश का, इसलिए
उसके बीस हज़ार \omterm{def:g10:universal-dna:gene}{जीनों} के बीच वह अकेला बाहरी \omterm{def:g10:universal-dna:gene}{जीन} भी बाक़ियों जैसा ही पढ़ा
जाता है।
\end{solution}
